% ============================================================
%  Rapport de projet — Systeme de gestion d'une pharmacie
%  Annee universitaire 2025/2026
% ============================================================
\documentclass[12pt,a4paper]{report}

% ── Encodage et langue ──────────────────────────────────────
\usepackage[utf8]{inputenc}
\usepackage[T1]{fontenc}
\usepackage[french]{babel}

% ── Mise en page ────────────────────────────────────────────
\usepackage[top=2.5cm,bottom=2.5cm,left=2.5cm,right=2.5cm]{geometry}
\usepackage{setspace}
\onehalfspacing
\setlength{\parindent}{0pt}
\setlength{\parskip}{6pt}

% ── Polices et couleurs ─────────────────────────────────────
\usepackage{lmodern}
\usepackage{xcolor}
\definecolor{bleutitre}{RGB}{0,70,130}

% ── Tableaux, figures, listes ───────────────────────────────
\usepackage{booktabs}
\usepackage{array}
\usepackage{longtable}
\usepackage{tabularx}
\usepackage{graphicx}
\usepackage{float}
\usepackage{enumitem}

% ── Code source ─────────────────────────────────────────────
\usepackage{listings}
\lstset{
  basicstyle=\ttfamily\small,
  breaklines=true,
  frame=single,
  backgroundcolor=\color{gray!8},
  keywordstyle=\color{blue!70!black},
  commentstyle=\color{green!50!black},
  stringstyle=\color{red!60!black},
  showstringspaces=false,
  tabsize=4
}

% ── Liens ───────────────────────────────────────────────────
\usepackage[hidelinks]{hyperref}

% ── Titres ──────────────────────────────────────────────────
\usepackage{titlesec}
\titleformat{\chapter}[hang]{\LARGE\bfseries\color{bleutitre}}{Chapitre \thechapter\;:\;}{0pt}{}
\titlespacing*{\chapter}{0pt}{-10pt}{18pt}
\titleformat{\section}[hang]{\large\bfseries\color{bleutitre}}{\thesection\quad}{0pt}{}
\titleformat{\subsection}[hang]{\normalsize\bfseries}{\thesubsection\quad}{0pt}{}

% ── En-tête / pied de page ─────────────────────────────────
\usepackage{fancyhdr}
\pagestyle{fancy}
\fancyhf{}
\fancyhead[L]{\small\textit{Système de gestion d'une pharmacie}}
\fancyhead[R]{\small\textit{2025/2026}}
\fancyfoot[C]{\thepage}
\renewcommand{\headrulewidth}{0.4pt}

% ============================================================
\begin{document}

% ════════════════════════════════════════════════════════════
%  PAGE DE GARDE
% ════════════════════════════════════════════════════════════
\begin{titlepage}
\centering
\vspace*{2cm}

{\large Université Hassan 1\textsuperscript{er} --- Faculté des Sciences et Techniques}\\[0.3cm]
{\large Département Informatique}\\[2.5cm]

{\Huge\bfseries\color{bleutitre} Système de gestion\\[0.3cm] d'une pharmacie}\\[1.5cm]

{\Large Projet de programmation Java}\\[0.4cm]
{\large Sujet 4}\\[2cm]

\begin{tabular}{rl}
\textbf{Encadrant}     & Pr. Said El Kafhali \\[6pt]
\textbf{Réalisé par}   & Oussama Tarchi \\
                        & Ali Zouine \\
                        & Taha Rajel \\
                        & Abdessamad Lyamani \\
\end{tabular}

\vfill
{\large Année universitaire 2025/2026}
\end{titlepage}

% ════════════════════════════════════════════════════════════
%  REMERCIEMENTS
% ════════════════════════════════════════════════════════════
\chapter*{Remerciements}
\addcontentsline{toc}{chapter}{Remerciements}

Nous tenons à exprimer notre gratitude à notre encadrant, le Professeur Said El Kafhali, pour le temps qu'il nous a consacré, ses orientations pertinentes et sa disponibilité tout au long de ce projet. Ses remarques nous ont permis de mieux structurer notre travail et de surmonter plusieurs difficultés techniques.

Nous remercions également les enseignants du département d'informatique pour la formation qu'ils nous ont dispensée, notamment en programmation orientée objet et en bases de données relationnelles, deux piliers sur lesquels repose ce projet.

Enfin, un remerciement particulier à nos familles et à nos camarades pour leur soutien et leurs encouragements constants.

% ════════════════════════════════════════════════════════════
%  TABLE DES MATIERES
% ════════════════════════════════════════════════════════════
\tableofcontents
\thispagestyle{empty}
\newpage

% ════════════════════════════════════════════════════════════
%  CHAPITRE 1 — INTRODUCTION GENERALE
% ════════════════════════════════════════════════════════════
\chapter{Introduction générale}

Une pharmacie, au-delà de son rôle médical, est un lieu où s'entrecroisent des flux de données variés : entrées de stock, sorties par vente, suivi de péremption, prescriptions médicales, relations avec les fournisseurs. Quand ces informations sont gérées manuellement --- sur des registres papier ou des fichiers tableurs dispersés --- le risque d'erreur devient considérable. Un médicament périmé peut rester en rayon, une rupture de stock peut passer inaperçue, et la traçabilité des ordonnances se perd dans l'accumulation des documents.

Le sujet qui nous a été confié dans le cadre de ce module vise justement à répondre à cette problématique : concevoir et développer une application desktop capable de centraliser l'ensemble des opérations courantes d'une pharmacie. Le pharmacien doit pouvoir, depuis une seule interface, enregistrer ses médicaments, gérer ses clients, saisir des ordonnances, effectuer des ventes, passer des commandes auprès de ses fournisseurs et surveiller l'état de son stock.

Pour réaliser ce système, nous avons choisi Java comme langage de développement, Swing pour l'interface graphique, MySQL pour le stockage des données et JDBC pour la communication entre l'application et la base de données. Le projet suit une architecture en couches --- Vue, Contrôleur, DAO, Modèle --- qui sépare clairement les responsabilités et facilite la maintenance du code.

Ce rapport retrace les étapes du projet : de l'analyse des besoins à la conception, puis de la réalisation technique aux tests, en passant par les difficultés rencontrées et les solutions adoptées. Nous terminerons par les limites de l'application et les perspectives d'amélioration envisagées.

% ════════════════════════════════════════════════════════════
%  CHAPITRE 2 — ANALYSE DU BESOIN
% ════════════════════════════════════════════════════════════
\chapter{Analyse du besoin}

\section{Présentation du sujet}

Le sujet demande la réalisation d'un système de gestion de pharmacie couvrant les opérations suivantes : gestion des médicaments et de leur stock, gestion des clients, saisie et vérification des ordonnances, enregistrement des ventes, suivi des commandes fournisseurs, historique des transactions et analyse des ventes.

L'objectif est de fournir un outil fonctionnel et intuitif permettant au pharmacien d'effectuer l'ensemble de ses tâches quotidiennes depuis une seule application, sans recourir à plusieurs outils distincts.

\section{Acteurs du système}

L'analyse du domaine fait apparaître quatre catégories d'acteurs qui interagissent, directement ou indirectement, avec le système :

\begin{table}[H]
\centering
\begin{tabularx}{\textwidth}{lX}
\toprule
\textbf{Acteur} & \textbf{Rôle} \\
\midrule
Pharmacien     & Utilisateur principal : enregistre les ventes, gère les médicaments, saisit les ordonnances, passe les commandes fournisseurs. \\
Responsable    & Consulte les historiques, les indicateurs du dashboard et les alertes de stock pour prendre les décisions d'approvisionnement. \\
Client         & Personne qui achète des médicaments ; ses coordonnées et son historique d'achats sont enregistrés dans le système. \\
Fournisseur    & Entreprise qui approvisionne la pharmacie en médicaments ; ses informations et ses commandes sont suivies. \\
\bottomrule
\end{tabularx}
\caption{Acteurs du système}
\end{table}

\section{Besoins fonctionnels}

Les besoins fonctionnels identifiés lors de l'analyse sont les suivants :

\begin{itemize}[nosep]
  \item Enregistrer et modifier les informations générales de la pharmacie.
  \item Ajouter, modifier, supprimer et rechercher des clients.
  \item Ajouter, modifier, supprimer et rechercher des médicaments.
  \item Suivre la quantité disponible en stock et détecter les seuils d'alerte.
  \item Détecter les médicaments périmés ou proches de la date de péremption.
  \item Gérer les fournisseurs (ajout, modification, suppression).
  \item Saisir les ordonnances et les associer à un client.
  \item Créer des ventes avec ou sans ordonnance.
  \item Mettre à jour automatiquement le stock après validation d'une vente.
  \item Créer des commandes fournisseurs et enregistrer leur réception.
  \item Augmenter le stock après réception d'une livraison.
  \item Générer automatiquement une facture texte après chaque vente validée.
  \item Consulter l'historique des ventes et des commandes.
  \item Analyser les ventes à travers un dashboard : par produit, par client, par période.
\end{itemize}

\section{Besoins non fonctionnels}

\begin{table}[H]
\centering
\begin{tabularx}{\textwidth}{lX}
\toprule
\textbf{Critère} & \textbf{Description} \\
\midrule
Simplicité       & L'interface doit être intuitive, exploitable sans formation préalable. \\
Fiabilité        & Les données doivent rester cohérentes après chaque opération : stock, ventes, commandes. \\
Rapidité         & Les opérations courantes doivent s'exécuter sans latence perceptible. \\
Persistance      & Toutes les données sont stockées dans MySQL et conservées entre les sessions. \\
\bottomrule
\end{tabularx}
\caption{Besoins non fonctionnels}
\end{table}

% ════════════════════════════════════════════════════════════
%  CHAPITRE 3 — CONCEPTION DU SYSTEME
% ════════════════════════════════════════════════════════════
\chapter{Conception du système}

\section{Architecture du projet}

Le projet adopte une architecture en quatre couches. Chaque couche a une responsabilité bien délimitée, ce qui facilite la maintenance et l'évolution du code.

\begin{table}[H]
\centering
\begin{tabularx}{\textwidth}{llX}
\toprule
\textbf{Couche} & \textbf{Package} & \textbf{Responsabilité} \\
\midrule
Vue         & \texttt{view}       & Affichage : formulaires, tableaux, boutons et onglets. \\
Contrôleur  & \texttt{controller} & Validation des données et application de la logique métier. \\
DAO         & \texttt{dao}        & Exécution des requêtes SQL et communication avec MySQL. \\
Modèle      & \texttt{model}      & Représentation des objets du domaine (Client, Médicament, etc.). \\
\bottomrule
\end{tabularx}
\caption{Couches de l'architecture}
\end{table}

Le flux de données suit un chemin linéaire : l'utilisateur interagit avec la vue, la vue délègue au contrôleur, le contrôleur appelle le DAO, et le DAO communique avec MySQL. Les objets modèle circulent entre ces couches pour transporter les données.

\begin{figure}[H]
\centering
\fbox{\parbox{0.85\textwidth}{\centering
\textbf{Utilisateur} $\longrightarrow$ \textbf{Vue (Swing)} $\longrightarrow$ \textbf{Contrôleur} $\longrightarrow$ \textbf{DAO (JDBC)} $\longrightarrow$ \textbf{MySQL}
}}
\caption{Flux de données dans l'architecture}
\end{figure}

\section{Structure des packages}

Le code source est organisé en cinq packages :

\begin{table}[H]
\centering
\begin{tabularx}{\textwidth}{lX}
\toprule
\textbf{Package} & \textbf{Contenu} \\
\midrule
\texttt{Main}       & Classe \texttt{Main} : point d'entrée de l'application. \\
\texttt{model}      & Neuf classes : \texttt{Client}, \texttt{Medicament}, \texttt{Fournisseur}, \texttt{Pharmacie}, \texttt{Ordonnance}, \texttt{Vente}, \texttt{LigneVente}, \texttt{Commande}, \texttt{LigneCommande}. \\
\texttt{dao}        & Un DAO par entité, plus \texttt{Utilitaire} (connexion) et \texttt{TestConnection}. \\
\texttt{controller} & Neuf contrôleurs dont \texttt{DashboardController} pour l'analyse des ventes. \\
\texttt{view}       & Fenêtre principale \texttt{GestionPharmacieView} et un panel par module, dont \texttt{DashboardPanel}. \\
\bottomrule
\end{tabularx}
\caption{Organisation des packages}
\end{table}

\section{Classes principales du modèle}

Le package \texttt{model} contient neuf classes qui représentent les entités du domaine :

\begin{table}[H]
\centering
\begin{tabularx}{\textwidth}{lX}
\toprule
\textbf{Classe} & \textbf{Description} \\
\midrule
\texttt{Client}         & Client avec nom, prénom, téléphone et adresse. \\
\texttt{Medicament}     & Médicament avec nom commercial, composition, forme, dosage, prix, quantité en stock, date de péremption et seuil d'alerte. \\
\texttt{Fournisseur}    & Fournisseur avec nom de société, contact et adresse. \\
\texttt{Pharmacie}      & Informations de l'établissement : nom, adresse, téléphone, email, responsable. \\
\texttt{Ordonnance}     & Ordonnance associée à un client, avec médecin, date et notes. \\
\texttt{Vente}          & Vente avec client, ordonnance optionnelle, date et liste de lignes. \\
\texttt{LigneVente}     & Détail d'un médicament vendu : identifiant vente, médicament, quantité, prix unitaire. \\
\texttt{Commande}       & Commande fournisseur avec date, statut et liste de lignes. \\
\texttt{LigneCommande}  & Détail d'un médicament commandé : identifiant commande, médicament, quantité, prix d'achat. \\
\bottomrule
\end{tabularx}
\caption{Classes du package \texttt{model}}
\end{table}

\section{Diagramme de cas d'utilisation}

Le diagramme ci-dessous identifie les interactions entre les acteurs et le système. Le pharmacien accède aux fonctionnalités opérationnelles, tandis que le responsable se concentre sur l'analyse et le pilotage.

\begin{figure}[H]
\centering
\fbox{\parbox{0.92\textwidth}{
\textbf{Pharmacien :}
\begin{itemize}[nosep,leftmargin=1.5em]
  \item Gérer les clients (ajout, modification, suppression)
  \item Gérer les médicaments
  \item Saisir une ordonnance
  \item Créer et valider une vente (avec ou sans ordonnance)
  \item Gérer le stock (entrées, sorties)
  \item Créer une commande fournisseur
  \item Réceptionner une livraison
  \item Générer une facture texte
\end{itemize}
\textbf{Responsable :}
\begin{itemize}[nosep,leftmargin=1.5em]
  \item Consulter les historiques de ventes et de commandes
  \item Analyser le stock faible et la péremption
  \item Calculer la valeur totale du stock
  \item Consulter le dashboard : ventes par produit, par client, chiffre d'affaires par période
  \item Planifier les approvisionnements
\end{itemize}
}}
\caption{Cas d'utilisation du système}
\end{figure}

\section{Diagramme de classes}

Le diagramme de classes représente les neuf entités du modèle et leurs associations :

\begin{itemize}[nosep]
  \item Un \texttt{Client} peut posséder plusieurs \texttt{Ordonnance} (1 $\rightarrow$ 0..*).
  \item Un \texttt{Client} peut effectuer plusieurs \texttt{Vente} (1 $\rightarrow$ 0..*).
  \item Une \texttt{Ordonnance} peut justifier une \texttt{Vente} (0..1 $\rightarrow$ 0..*).
  \item Une \texttt{Vente} contient une ou plusieurs \texttt{LigneVente} (1 $\rightarrow$ 1..*).
  \item Une \texttt{LigneVente} concerne un \texttt{Medicament} (0..* $\rightarrow$ 1).
  \item Un \texttt{Fournisseur} peut fournir plusieurs \texttt{Commande} (1 $\rightarrow$ 0..*).
  \item Une \texttt{Commande} contient une ou plusieurs \texttt{LigneCommande} (1 $\rightarrow$ 1..*).
  \item Une \texttt{LigneCommande} approvisionne un \texttt{Medicament} (0..* $\rightarrow$ 1).
\end{itemize}

\begin{figure}[H]
\centering
\fbox{\parbox{0.92\textwidth}{\small
\begin{verbatim}
 Client ──(1)──< possede >──(0..*)── Ordonnance
 Client ──(1)──< effectue >──(0..*)── Vente
 Ordonnance ──(0..1)──< justifie >──(0..*)── Vente
 Vente ──(1)──*< contient >*(1..*)── LigneVente
 LigneVente ──(0..*)──< concerne >──(1)── Medicament
 Fournisseur ──(1)──< fournit >──(0..*)── Commande
 Commande ──(1)──*< contient >*(1..*)── LigneCommande
 LigneCommande ──(0..*)──< approvisionne >──(1)── Medicament
\end{verbatim}
}}
\caption{Relations entre les classes du modèle}
\end{figure}

\section{Diagramme entité-relation}

Le diagramme entité-relation correspond directement à la structure de la base de données. Chaque entité devient une table, chaque association se traduit par une clé étrangère.

\begin{figure}[H]
\centering
\fbox{\parbox{0.92\textwidth}{\small
\begin{verbatim}
 CLIENT(id PK, nom, prenom, telephone, adresse)
 MEDICAMENT(id PK, nomCommercial, composition, forme,
            dosage, prix, quantiteStock, datePeremption,
            seuilAlerte)
 FOURNISSEUR(id PK, nomSociete, contact, adresse)
 PHARMACIE(id PK, nom, adresse, telephone, email, responsable)
 ORDONNANCE(id PK, idClient FK->CLIENT, medecin,
            dateOrdonnance, notes)
 VENTE(id PK, idClient FK->CLIENT,
       idOrdonnance FK->ORDONNANCE, dateVente)
 LIGNE_VENTE(idVente FK->VENTE, idMedicament FK->MEDICAMENT,
             quantite, prixUnitaire)
 COMMANDE(id PK, idFournisseur FK->FOURNISSEUR,
          dateCommande, statut)
 LIGNE_COMMANDE(idCommande FK->COMMANDE,
                idMedicament FK->MEDICAMENT,
                quantiteCommandee, prixAchat)
\end{verbatim}
}}
\caption{Schéma entité-relation de la base \texttt{gestion\_pharmacie}}
\end{figure}

% ════════════════════════════════════════════════════════════
%  CHAPITRE 4 — BASE DE DONNEES
% ════════════════════════════════════════════════════════════
\chapter{Base de données}

\section{Présentation de la base}

La base de données, nommée \texttt{gestion\_pharmacie}, fonctionne sous MySQL. Elle contient neuf tables correspondant aux neuf entités du modèle. Le script de création définit les tables, les colonnes, les types de données et les contraintes d'intégrité.

\section{Tables et relations}

\begin{table}[H]
\centering
\begin{tabularx}{\textwidth}{lX}
\toprule
\textbf{Table} & \textbf{Contenu et relations} \\
\midrule
\texttt{Client}         & Clients de la pharmacie. \\
\texttt{Medicament}     & Médicaments avec prix, stock et péremption. \\
\texttt{Fournisseur}    & Fournisseurs de médicaments. \\
\texttt{Pharmacie}      & Fiche de l'établissement. \\
\texttt{Ordonnance}     & Ordonnances (\texttt{idClient} $\rightarrow$ \texttt{Client.id}). \\
\texttt{Vente}          & Ventes (\texttt{idClient} $\rightarrow$ \texttt{Client.id}, \texttt{idOrdonnance} $\rightarrow$ \texttt{Ordonnance.id}). \\
\texttt{LigneVente}     & Détail des ventes (\texttt{idVente} $\rightarrow$ \texttt{Vente.id}, \texttt{idMedicament} $\rightarrow$ \texttt{Medicament.id}). \\
\texttt{Commande}       & Commandes fournisseurs (\texttt{idFournisseur} $\rightarrow$ \texttt{Fournisseur.id}). \\
\texttt{LigneCommande}  & Détail des commandes (\texttt{idCommande} $\rightarrow$ \texttt{Commande.id}, \texttt{idMedicament} $\rightarrow$ \texttt{Medicament.id}). \\
\bottomrule
\end{tabularx}
\caption{Tables de la base de données}
\end{table}

\section{Clés primaires et étrangères}

Chaque table possède une clé primaire \texttt{id} auto-incrémentée, sauf \texttt{LigneVente} et \texttt{LigneCommande} qui utilisent des clés composées (couple \texttt{idVente}/\texttt{idMedicament} et \texttt{idCommande}/\texttt{idMedicament}). Les clés étrangères assurent l'intégrité référentielle : il est impossible d'insérer une vente référençant un client inexistant, ou une ligne de vente pointant vers un médicament supprimé.

Cette structuration garantit la cohérence des données à chaque opération. Par exemple, la suppression d'un client n'est possible que si aucune ordonnance ni aucune vente ne lui est rattachée.

\section{Correspondance noms SQL et DAO}

Au cours du développement, nous avons rencontré des incohérences entre les noms de colonnes définis dans les tables MySQL et les noms utilisés dans les requêtes SQL des classes DAO. Un attribut nommé \texttt{nomSociete} dans la table \texttt{Fournisseur} pouvait, par inattention, être écrit différemment dans le code Java. Ce décalage provoquait des exceptions \texttt{SQLException} empêchant l'affichage ou l'enregistrement des données.

La solution a consisté à comparer systématiquement la structure de chaque table avec les requêtes du DAO correspondant, puis à aligner les noms de colonnes. Ce travail minutieux a été réalisé sur l'ensemble des sept DAO du projet.

% ════════════════════════════════════════════════════════════
%  CHAPITRE 5 — REALISATION TECHNIQUE
% ════════════════════════════════════════════════════════════
\chapter{Réalisation technique}

\section{Interface graphique : Java Swing}

L'interface a été construite avec Java Swing. La fenêtre principale, \texttt{GestionPharmacieView}, utilise un composant \texttt{JTabbedPane} qui organise les modules en onglets. L'utilisateur navigue entre les fonctionnalités en cliquant simplement sur l'onglet concerné.

\begin{table}[H]
\centering
\begin{tabularx}{\textwidth}{lX}
\toprule
\textbf{Onglet} & \textbf{Fonction} \\
\midrule
Pharmacie    & Formulaire d'enregistrement des informations de l'établissement. \\
Clients      & Ajout, modification, suppression et recherche de clients. \\
Médicaments  & Ajout, modification, suppression et recherche de médicaments. \\
Fournisseurs & Gestion des fournisseurs. \\
Ordonnances  & Saisie et consultation des ordonnances. \\
Ventes       & Création de ventes, panier, validation, ordonnance et historique. \\
Stock        & Suivi du stock, alertes de stock faible, péremption et valeur du stock. \\
Dashboard    & Tableaux d'analyse : ventes par produit, par client, produits les plus demandés et chiffre d'affaires par période. \\
\bottomrule
\end{tabularx}
\caption{Onglets de l'interface principale}
\end{table}

La fenêtre mesure $1000 \times 650$ pixels par défaut, est centrée à l'écran au démarrage et se ferme avec \texttt{EXIT\_ON\_CLOSE}.

\textbf{Remarque :} Le module \texttt{CommandePanel} pour les commandes fournisseurs a été entièrement développé et testé. Il est temporairement désactivé dans la fenêtre principale (la ligne d'ajout de l'onglet est commentée dans \texttt{GestionPharmacieView}) pour des raisons d'intégration en cours de finalisation, mais le code complet est présent et fonctionnel.

\section{Accès aux données : JDBC}

La communication avec MySQL passe par JDBC avec le connecteur \texttt{mysql-connector-j-9.7.0.jar}. La classe \texttt{Utilitaire} centralise la gestion de la connexion :

\begin{lstlisting}[language=Java,caption={Classe Utilitaire — connexion MySQL}]
public class Utilitaire {
    private static final String URL =
        "jdbc:mysql://localhost:3306/gestion_pharmacie";
    private static final String USER = "root";
    private static final String PASSWORD = "";
    private static Connection conn = null;

    public static Connection getConnection() {
        if (conn == null || conn.isClosed()) {
            Class.forName("com.mysql.cj.jdbc.Driver");
            conn = DriverManager.getConnection(URL, USER, PASSWORD);
        }
        return conn;
    }

    public static void closeConnection() {
        if (conn != null && !conn.isClosed()) conn.close();
    }
}
\end{lstlisting}

La classe maintient une connexion unique (\emph{singleton}) réutilisée par tous les DAO. La méthode \texttt{closeConnection()} permet de libérer la ressource à la fin de l'utilisation.

\section{Couche DAO}

Chaque entité dispose d'un DAO dédié. Les DAO utilisent des \texttt{PreparedStatement} pour exécuter les requêtes SQL, ce qui protège contre les injections SQL. Les opérations standard sont :

\begin{itemize}[nosep]
  \item \texttt{ajouter()} : insertion d'un nouvel enregistrement.
  \item \texttt{modifier()} : mise à jour d'un enregistrement existant.
  \item \texttt{supprimer()} : suppression d'un enregistrement.
  \item \texttt{getAll()} : récupération de la liste complète.
  \item \texttt{getById()} : recherche par identifiant.
\end{itemize}

Certains DAO possèdent des méthodes supplémentaires. Le \texttt{VenteDao} gère l'enregistrement des lignes de vente (\texttt{enregistrerLigne}), la mise à jour du stock via une requête \texttt{UPDATE} décrémentant la quantité disponible, et la modification de l'ordonnance associée (\texttt{modifierOrdo}). Le \texttt{CommandeDao} gère le changement de statut et la réception des livraisons.

\section{Couche Contrôleur}

Les contrôleurs encapsulent la logique métier et évitent de placer des règles de gestion dans les vues ou les DAO. Voici quelques exemples concrets :

\begin{itemize}[nosep]
  \item \texttt{VenteController} vérifie que le stock est suffisant avant de valider une vente, s'assure que l'ordonnance correspond bien au client, et génère la facture texte après validation.
  \item \texttt{StockController} calcule la valeur totale du stock et génère les alertes de stock faible et de péremption.
  \item \texttt{CommandeController} gère le changement de statut d'une commande (En attente $\rightarrow$ Validée $\rightarrow$ Livrée) et la mise à jour du stock à la réception.
  \item \texttt{DashboardController} agrège les données de ventes pour produire des statistiques par produit, par client et par période.
\end{itemize}

\section{Génération des factures texte}

Après la validation d'une vente, le \texttt{VenteController} génère automatiquement une facture sous forme de fichier texte. La méthode \texttt{genererEtEnregistrerFacture()} construit le contenu de la facture (numéro, client, date, détail des lignes et total) puis l'écrit dans un fichier nommé \texttt{facture\_<id>.txt} dans le dossier \texttt{factures/}. Si ce dossier n'existe pas, il est créé automatiquement.

La facture contient :
\begin{itemize}[nosep]
  \item Le numéro de la vente.
  \item L'identifiant du client.
  \item La date de vente.
  \item La mention de l'ordonnance (numéro ou \og Sans ordonnance \fg).
  \item Le détail de chaque ligne : médicament, quantité, prix unitaire et sous-total.
  \item Le total à payer en dirhams.
\end{itemize}

\section{Point d'entrée}

La classe \texttt{Main.Main} lance l'application. Elle utilise \texttt{SwingUtilities.invokeLater()} pour respecter le modèle de threading de Swing :

\begin{lstlisting}[language=Java,caption={Classe Main}]
public class Main {
    public static void main(String[] args) {
        SwingUtilities.invokeLater(new Runnable() {
            public void run() {
                GestionPharmacieView fenetre =
                    new GestionPharmacieView();
                fenetre.setVisible(true);
            }
        });
    }
}
\end{lstlisting}

% ════════════════════════════════════════════════════════════
%  CHAPITRE 6 — FONCTIONNEMENT DE L'APPLICATION
% ════════════════════════════════════════════════════════════
\chapter{Fonctionnement de l'application}

\section{Ajout d'un médicament}

L'utilisateur ouvre l'onglet \og Médicaments \fg, remplit le formulaire avec le nom commercial, la composition, la forme, le dosage, le prix, la quantité en stock, la date de péremption et le seuil d'alerte, puis clique sur \og Ajouter \fg. Le \texttt{MedicamentController} valide les champs et appelle le DAO qui exécute un \texttt{INSERT} dans la table \texttt{Medicament}. Le tableau se rafraîchit automatiquement pour afficher le nouveau produit.

\section{Création d'un client}

Depuis l'onglet \og Clients \fg, le pharmacien saisit le nom, le prénom, le téléphone et l'adresse. Le \texttt{ClientController} vérifie que les champs obligatoires sont remplis, puis le \texttt{ClientDao} insère l'enregistrement. Le client est immédiatement visible dans la liste et disponible dans les menus déroulants des autres modules.

\section{Création d'une ordonnance}

L'onglet \og Ordonnances \fg{} permet de saisir l'identifiant du client, le nom du médecin, la date et des notes éventuelles. L'\texttt{OrdonnanceController} vérifie que le client existe dans la base et que le nom du médecin est renseigné, puis l'\texttt{OrdonnanceDao} insère l'ordonnance.

\section{Vente avec ou sans ordonnance}

Le scénario de vente se déroule en plusieurs étapes :

\begin{enumerate}[nosep]
  \item Le pharmacien sélectionne un client dans la liste déroulante (ou choisit \og Client anonyme \fg{} pour une vente sans identification).
  \item Il sélectionne un médicament et choisit la quantité souhaitée.
  \item Il clique sur \og Ajouter ligne \fg. Si c'est la première ligne, une vente est créée automatiquement dans la base de données. L'identifiant de cette vente est conservé dans la variable \texttt{idVenteEnCours}.
  \item Il peut ajouter d'autres lignes au panier. Le sous-total et le total se mettent à jour en temps réel.
  \item S'il coche la case \og Ordonnance présente \fg, des champs supplémentaires apparaissent pour saisir le médecin, la date et les notes.
  \item Il clique sur \og Valider \fg. Le contrôleur vérifie que le stock est suffisant pour chaque ligne. Si une ordonnance a été saisie, elle est enregistrée dans la table \texttt{Ordonnance} puis associée à la vente via \texttt{modifierOrdo}.
  \item Le stock de chaque médicament est décrémenté de la quantité vendue.
  \item Une facture texte est automatiquement générée dans le dossier \texttt{factures/}.
  \item La vente apparaît dans l'historique en bas de l'écran.
\end{enumerate}

\section{Génération automatique de facture}

Immédiatement après la validation de la vente, la méthode \texttt{genererEtEnregistrerFacture()} est appelée. Elle produit un fichier \texttt{facture\_<id>.txt} contenant l'en-tête, le détail des achats et le total. Ce fichier est stocké sur le disque dans le dossier \texttt{factures/}, créé automatiquement si nécessaire.

\section{Commande fournisseur et réception}

Le pharmacien crée une commande en sélectionnant un fournisseur, un médicament, une quantité et un prix d'achat. La commande est enregistrée avec le statut \og En attente \fg. Il peut ensuite la valider (statut \og Validée \fg), puis la réceptionner (statut \og Livrée \fg). Lors de la réception, le stock de chaque médicament commandé est automatiquement augmenté de la quantité livrée.

Le \texttt{CommandePanel} affiche également un historique des commandes avec le fournisseur, la date, le statut et le total.

\section{Suivi du stock et alertes}

L'onglet \og Stock \fg{} offre une vue d'ensemble sur les médicaments disponibles. Le \texttt{StockController} propose les fonctions suivantes :

\begin{itemize}[nosep]
  \item \textbf{Stock faible} : affiche les médicaments dont la quantité est inférieure ou égale au seuil d'alerte.
  \item \textbf{Produits périmés} : affiche les médicaments dont la date de péremption est dépassée.
  \item \textbf{Péremption proche} : affiche les médicaments qui expireront dans les 30 prochains jours.
  \item \textbf{Valeur du stock} : calcule la somme $\text{prix} \times \text{quantité}$ pour l'ensemble des médicaments.
\end{itemize}

\section{Analyse avec le dashboard}

Le module Dashboard, accessible depuis l'onglet éponyme, centralise plusieurs indicateurs construits à partir des ventes enregistrées :

\begin{itemize}[nosep]
  \item \textbf{Ventes par produit} : un tableau affiche, pour chaque médicament, la quantité totale vendue.
  \item \textbf{Ventes par client} : un tableau affiche, pour chaque client, le montant total de ses achats en dirhams.
  \item \textbf{Produits les plus demandés} : les médicaments sont classés par quantité vendue décroissante.
  \item \textbf{Chiffre d'affaires par période} : l'utilisateur saisit une date de début et une date de fin au format \texttt{yyyy-mm-dd}, puis clique sur \og Calculer CA \fg. Le résultat s'affiche dans un champ dédié.
  \item Le bouton \og Actualiser tout \fg{} recharge l'ensemble des statistiques depuis la base.
\end{itemize}

Le \texttt{DashboardController} exploite les données de \texttt{VenteController}, \texttt{ClientController} et \texttt{MedicamentController} pour agréger les résultats. Le \texttt{DashboardPanel} utilise des sous-onglets \texttt{JTabbedPane} pour organiser les différentes vues d'analyse.

% ════════════════════════════════════════════════════════════
%  CHAPITRE 7 — PROBLEMES RENCONTRES
% ════════════════════════════════════════════════════════════
\chapter{Problèmes rencontrés}

Ce chapitre détaille les principaux problèmes rencontrés pendant le développement, avec pour chacun la cause identifiée, l'impact sur l'application et la solution mise en place.

\section{Incohérence entre noms SQL et DAO}

\textbf{Cause :} Certains noms de colonnes dans les tables MySQL ne correspondaient pas exactement aux noms utilisés dans les requêtes des classes DAO. Par exemple, \texttt{nomSociete} dans la table \texttt{Fournisseur} pouvait être écrit différemment dans le code Java.

\textbf{Impact :} L'exécution des requêtes provoquait des exceptions \texttt{SQLException}, empêchant l'affichage ou l'enregistrement des données.

\textbf{Solution :} Nous avons comparé méthodiquement la structure de chaque table MySQL avec les requêtes de chaque DAO, puis aligné les noms pour qu'ils correspondent exactement.

\section{Intégration de l'ordonnance dans l'interface de vente}

\textbf{Cause :} L'ordonnance doit être associée à la fois au client et à la vente. Sa saisie depuis l'interface de vente nécessitait des champs conditionnels et une gestion précise du moment de l'enregistrement.

\textbf{Impact :} L'ordonnance risquait d'être créée avant la vente, ce qui posait un problème de cohérence.

\textbf{Solution :} Nous avons ajouté une case à cocher \og Ordonnance présente \fg{} dans le \texttt{VentePanel}. Les champs (médecin, date, notes) ne s'affichent que lorsque la case est cochée. L'ordonnance n'est enregistrée qu'au moment de la validation de la vente, jamais avant.

\section{Ajout d'une ligne avant création de la vente}

\textbf{Cause :} Une ligne de vente nécessite un \texttt{idVente} pour être insérée dans la table \texttt{LigneVente}. Or, la vente n'existait pas encore lorsque l'utilisateur ajoutait la première ligne.

\textbf{Impact :} L'insertion échouait car la clé étrangère \texttt{idVente} était invalide.

\textbf{Solution :} Au premier ajout de ligne, la vue crée automatiquement une vente dans la base et conserve son identifiant dans \texttt{idVenteEnCours}. Les lignes suivantes réutilisent cet identifiant.

\section{Association tardive de l'ordonnance à la vente}

\textbf{Cause :} L'ordonnance est créée au moment de la validation, alors que la vente existe déjà depuis le premier ajout de ligne. Il fallait donc pouvoir modifier la vente après sa création initiale.

\textbf{Impact :} La vente restait sans ordonnance même lorsque le pharmacien en avait saisi une.

\textbf{Solution :} Ajout de la méthode \texttt{modifierOrdo()} dans \texttt{VenteController} et \texttt{VenteDao} pour mettre à jour le champ \texttt{idOrdonnance} dans la table \texttt{Vente} après l'enregistrement de l'ordonnance.

\section{Génération et stockage des factures texte}

\textbf{Cause :} La facture devait être générée sous forme de fichier texte et conservée physiquement sur le disque après validation de la vente.

\textbf{Impact :} Sans cette fonctionnalité, aucune trace imprimable de la transaction n'était conservée.

\textbf{Solution :} Ajout de la méthode \texttt{genererEtEnregistrerFacture()} dans \texttt{VenteController}. Elle crée le dossier \texttt{factures/} si nécessaire et y écrit un fichier \texttt{facture\_<id>.txt} contenant le détail complet de la vente.

\section{Ajout du dashboard analytique}

\textbf{Cause :} Le sujet demandait une partie analyse et reporting. L'ajout du module Dashboard a nécessité la création d'un contrôleur (\texttt{DashboardController}) et d'un panel (\texttt{DashboardPanel}) supplémentaires, ainsi que la coordination avec les contrôleurs existants.

\textbf{Impact :} L'intégration a demandé de s'assurer que les données de ventes, de clients et de médicaments étaient correctement agrégées sans duplication.

\textbf{Solution :} Le \texttt{DashboardController} réutilise les contrôleurs existants (\texttt{VenteController}, \texttt{ClientController}, \texttt{MedicamentController}) plutôt que d'accéder directement aux DAO, ce qui garantit la cohérence des calculs.

\section{Encodage UTF-8 sous Windows}

\textbf{Cause :} Le compilateur \texttt{javac} utilisait par défaut l'encodage Windows (CP1252), ce qui provoquait des erreurs lorsque le code source contenait des caractères accentués.

\textbf{Impact :} La compilation échouait avec des erreurs de caractères non reconnus.

\textbf{Solution :} Ajout de l'option \texttt{-encoding UTF-8} lors de la compilation en ligne de commande.

\section{Avertissement sur le nom du module Java}

\textbf{Cause :} Le nom initial du module Java se terminait par un chiffre (\texttt{gestion\_pharmacie2}), ce qui déclenchait un avertissement du compilateur.

\textbf{Impact :} Un message d'avertissement apparaissait à chaque compilation, sans bloquer l'exécution.

\textbf{Solution :} Le module a été renommé en \texttt{gestion.pharmacie}, conforme aux conventions de nommage Java.

% ════════════════════════════════════════════════════════════
%  CHAPITRE 8 — TESTS ET VALIDATION
% ════════════════════════════════════════════════════════════
\chapter{Tests et validation}

Les tests ont été réalisés manuellement en suivant les scénarios d'utilisation décrits dans le chapitre précédent. Chaque fonctionnalité a été vérifiée après son implémentation.

\begin{table}[H]
\centering
\begin{tabularx}{\textwidth}{Xcc}
\toprule
\textbf{Test effectué} & \textbf{Méthode} & \textbf{Résultat} \\
\midrule
Compilation du projet avec \texttt{javac -encoding UTF-8} & Terminal & OK \\
Connexion MySQL via \texttt{TestConnection} & Terminal & OK \\
Ajout, modification et suppression d'un client & Interface & OK \\
Ajout, modification et suppression d'un médicament & Interface & OK \\
Recherche d'un médicament par nom & Interface & OK \\
Ajout et suppression d'un fournisseur & Interface & OK \\
Création d'une ordonnance & Interface & OK \\
Création d'une vente sans ordonnance & Interface & OK \\
Création d'une vente avec ordonnance & Interface & OK \\
Mise à jour du stock après vente & Vérification en base & OK \\
Génération du fichier facture dans \texttt{factures/} & Système de fichiers & OK \\
Création d'une commande fournisseur & Interface & OK \\
Réception d'une livraison et mise à jour du stock & Vérification en base & OK \\
Détection du stock faible & Interface & OK \\
Détection des produits périmés & Interface & OK \\
Détection de la péremption proche & Interface & OK \\
Calcul de la valeur totale du stock & Interface & OK \\
Dashboard : ventes par produit & Interface & OK \\
Dashboard : ventes par client & Interface & OK \\
Dashboard : chiffre d'affaires par période & Interface & OK \\
Dashboard : produits les plus demandés & Interface & OK \\
\bottomrule
\end{tabularx}
\caption{Résultats des tests}
\end{table}

La classe \texttt{TestConnection} dans le package \texttt{dao} permet de vérifier rapidement que la connexion à MySQL fonctionne. En cas de succès, elle affiche le message \og Connexion MySQL réussie ! \fg{} dans la console.

Les tests du dashboard ont porté sur la cohérence des chiffres affichés par rapport aux données présentes dans la base. Le chiffre d'affaires par période a été validé en comparant le résultat affiché avec un calcul manuel effectué sur les lignes de vente de la période concernée.

% ════════════════════════════════════════════════════════════
%  CHAPITRE 9 — LIMITES ET PERSPECTIVES
% ════════════════════════════════════════════════════════════
\chapter{Limites et perspectives}

L'application remplit les objectifs fixés par le sujet, mais plusieurs améliorations sont envisageables pour la rendre plus complète et plus robuste :

\begin{enumerate}[nosep]
  \item \textbf{Graphiques pour le dashboard} : intégrer des histogrammes et des courbes (via une bibliothèque comme JFreeChart) pour visualiser l'évolution des ventes et du stock, plutôt que de s'appuyer uniquement sur des tableaux textuels.

  \item \textbf{Export PDF des factures} : permettre l'exportation des factures en format PDF pour l'impression et l'archivage professionnel, à l'aide d'une bibliothèque comme iText.

  \item \textbf{Authentification} : ajouter un écran de connexion avec un système de mots de passe pour restreindre l'accès à l'application et protéger les données sensibles.

  \item \textbf{Statistiques graphiques avancées} : proposer des analyses plus fines, comme l'évolution mensuelle du chiffre d'affaires ou les tendances de consommation par catégorie de médicament.

  \item \textbf{Tests automatisés} : mettre en place des tests unitaires avec JUnit pour vérifier automatiquement le comportement des contrôleurs et des DAO, réduisant ainsi le risque de régression.

  \item \textbf{Fichier de configuration} : externaliser les paramètres de connexion MySQL dans un fichier \texttt{.properties} au lieu de les coder en dur dans la classe \texttt{Utilitaire}, ce qui faciliterait le déploiement sur différents environnements.
\end{enumerate}

% ════════════════════════════════════════════════════════════
%  CHAPITRE 10 — CONCLUSION
% ════════════════════════════════════════════════════════════
\chapter{Conclusion}

Ce projet nous a permis de concevoir et de développer une application complète de gestion de pharmacie en Java. Le système couvre les opérations principales demandées dans le sujet : gestion des médicaments, suivi du stock avec alertes, gestion des clients, saisie des ordonnances, ventes avec ou sans prescription, commandes fournisseurs avec réception et mise à jour du stock, génération automatique de factures texte et analyse des ventes via un dashboard dédié.

L'architecture en couches que nous avons adoptée --- Vue, Contrôleur, DAO, Modèle --- structure le projet de manière lisible et maintenable. Les vues gèrent l'affichage et l'interaction avec l'utilisateur, les contrôleurs appliquent les règles métier, les DAO communiquent avec MySQL et les modèles représentent les données du domaine. Cette séparation nous a permis de travailler en parallèle sur différentes parties du code sans créer de conflits.

Les difficultés rencontrées au fil du développement --- incohérences entre noms SQL et DAO, gestion de l'ordonnance dans le processus de vente, création de la vente avant l'ajout des lignes, encodage Windows --- nous ont appris à diagnostiquer les erreurs, à lire les messages d'exception et à corriger le code de façon méthodique. L'ajout du module Dashboard et de la génération de factures a enrichi notre compréhension de l'intégration entre modules.

Sur le plan des compétences, ce projet nous a permis de consolider nos acquis en Java, Swing, JDBC et MySQL. Nous avons compris concrètement comment structurer un projet en couches, comment connecter une application desktop à une base de données relationnelle, et comment gérer les cas particuliers qui émergent lors de l'intégration de modules interdépendants.

L'application constitue une base solide, fonctionnelle et cohérente, qui pourrait être enrichie par de l'authentification, de l'export PDF, des graphiques et des tests automatisés pour répondre à des besoins plus avancés dans un contexte professionnel.

% ════════════════════════════════════════════════════════════
\end{document}
